\documentclass[11pt]{article}

\usepackage[T1]{fontenc}
\usepackage[utf8]{inputenc}
\usepackage{lmodern}
\usepackage{microtype}
\usepackage[margin=1in]{geometry}
\usepackage{amsmath,amssymb,amsthm,mathtools}
\usepackage{booktabs}
\usepackage{array}
\usepackage{enumitem}
\usepackage{xcolor}
\usepackage{listings}
\usepackage{fancyhdr}
\usepackage{hyperref}
\usepackage{url}

\hypersetup{
  colorlinks=true,
  linkcolor=blue!50!black,
  urlcolor=blue!50!black,
  citecolor=blue!50!black,
  pdfauthor={OpenAI research assistant},
  pdftitle={Seeded Reflection Power: A Proposed BusyLean Contender 7}
}

\pagestyle{fancy}
\fancyhf{}
\lhead{Seeded Reflection Power}
\rhead{BusyLean Contender 7 proposal}
\cfoot{\thepage}
\setlength{\headheight}{14pt}
\setlength{\emergencystretch}{2em}

\lstdefinelanguage{LeanFour}{
  keywords={import,namespace,end,open,abbrev,def,theorem,inductive,structure,where,
    match,with,fun,let,in,by,have,show,exact,rw,simpa,intro,partial,unsafe,
    noncomputable,termination_by},
  sensitive=true,
  morecomment=[l]{--},
  morecomment=[s]{/-}{-/},
  morestring=[b]"
}
\lstset{
  language=LeanFour,
  basicstyle=\ttfamily\footnotesize,
  keywordstyle=\bfseries\color{blue!45!black},
  commentstyle=\itshape\color{black!58},
  stringstyle=\color{green!35!black},
  frame=single,
  framerule=0.25pt,
  rulecolor=\color{black!25},
  xleftmargin=0.7em,
  xrightmargin=0.7em,
  columns=fullflexible,
  keepspaces=true,
  breaklines=true,
  breakatwhitespace=false,
  showstringspaces=false,
  tabsize=2,
  literate={→}{{$\to$}}1 {←}{{$\leftarrow$}}1 {∀}{{$\forall$}}1
           {∃}{{$\exists$}}1 {≤}{{$\le$}}1 {≥}{{$\ge$}}1
           {ℕ}{{$\mathbb N$}}1 {λ}{{$\lambda$}}1 {·}{{$\cdot$}}1
           {⟨}{{$\langle$}}1 {⟩}{{$\rangle$}}1
}

\newtheorem{theorem}{Theorem}[section]
\newtheorem{lemma}[theorem]{Lemma}
\newtheorem{proposition}[theorem]{Proposition}
\newtheorem{corollary}[theorem]{Corollary}
\theoremstyle{definition}
\newtheorem{definition}[theorem]{Definition}
\newtheorem{construction}[theorem]{Construction}
\theoremstyle{remark}
\newtheorem{remark}[theorem]{Remark}
\newtheorem{caution}[theorem]{Caution}

\newcommand{\Nat}{\mathbb N}
\newcommand{\Fam}{\mathcal F}
\newcommand{\Tree}{\mathcal W}
\newcommand{\Lift}{\mathcal L}
\newcommand{\Idx}{I}
\newcommand{\Stage}{\mathsf S}
\newcommand{\Pfive}{P}
\newcommand{\Height}{H}
\newcommand{\lean}[1]{\texttt{\detokenize{#1}}}
\newcommand{\repo}{\textsc{BusyLean}}
\newcommand{\code}[1]{\texttt{\detokenize{#1}}}

\title{\textbf{Seeded Reflection Power}\\[0.35em]
  \Large A Proposed Contender 7 for BusyLean}
\author{Research package prepared for Vladimir Reshetnikov}
\date{October 12, 2026}

\begin{document}
\maketitle

\begin{abstract}
We propose a new closed, constructive Lean natural number that strictly exceeds
\repo's current Contender~6 without using the incumbent value as an iteration
counter.  The key observation is that the existing forest-reflection scheduler
has a reusable higher-order type: it maps an oracle family and a W-tree control
object to a natural-valued function.  We define an order-raising operator that
reflects an \emph{entire scheduler} through the repository's universal W-tree
index.  Iterating this operator yields a strict hierarchy of schedulers.  The
final number uses a promotion height equal to the P5 seed value at~42, with a
small successor guard, and evaluates the resulting scheduler at the same
universal index and input~42.

The supplied Lean development is designed to prove
\[
  \texttt{BusyLean.contender\_6}
  < \texttt{BusyLean.contender\_7},
\]
as well as the inherited comparisons with the exact Contender~5 port and P5.
All recursion is structural or primitive, and the defining closure is intended
to avoid every named incumbent value.  The archive includes Lean source, a
kernel-side closure audit, a patch, a pinned CI workflow, executable finite
models, static audits, this article, and a validation ledger.  The present
environment did not contain a Lean executable, so the package does not claim
kernel verification; it records that limitation explicitly and supplies the
exact commands needed to complete the check under Lean~4.34.0.
\end{abstract}

\noindent\fbox{\parbox{0.965\linewidth}{
\textbf{Verification status.} The mathematical construction and proof term have
been written against the APIs in \repo{} commit
\code{86b360648163f0249eab79539438b6a83dfe8215}.  Static audits and executable
finite analogues pass.  The Lean files have \emph{not} been kernel-checked in
this environment.  No statement below should be read as a claim of completed
Lean verification until \code{lake build} and the supplied audit file both
succeed on the pinned toolchain.
}}

\tableofcontents

%=============================================================================
\section{Executive summary}

The incumbent has the form
\begin{equation}\label{eq:c6}
 C_6=\operatorname{stage}(\operatorname{p5Family},\Idx,42),
\end{equation}
where \(\Idx\) is a universal W-tree assembled from all trees denoted by the
closed programs of P5's W-language.  The repository informally calibrates this
near \(f_{\Lambda^2}(42)\), with
\(\Lambda\approx\psi(\Omega_\omega)\).

The new construction promotes the scheduler itself.  If \(S\) is a scheduler,
its lift \(\Lift S\) recursively traverses a second W-tree.  At a leaf it runs
\(S\) through \(\Idx\).  At a node it recursively lifts every child and exposes
those values to \(S\) as a new oracle family, again running \(S\) through
\(\Idx\).  Let
\[
 S_0=\operatorname{stage},\qquad S_{k+1}=\Lift S_k.
\]
Let
\[
 \Pfive=\operatorname{p5Family}(0,42),\qquad
 \Height=\operatorname{succ}(42+\Pfive).
\]
The proposed number is
\begin{equation}\label{eq:c7}
 C_7=S_{\Height}(\operatorname{p5Family},\Idx,42).
\end{equation}

The formal comparison rests on a single invariant:
\[
 \operatorname{Dom}(S):\iff
 \forall a\,\forall i\le n,
 \quad a(i,n)<S(a,\Idx,n).
\]
The existing scheduler has this property.  If \(S\) has it, then
\(\Lift S\) has it and, at \(\Idx\), strictly exceeds \(S\).  Therefore
\[
 S_0(a,\Idx,n)<S_1(a,\Idx,n)<S_2(a,\Idx,n)<\cdots.
\]
Since \(C_6=S_0(\operatorname{p5Family},\Idx,42)\) and \(\Height>0\), the
required strict inequality follows.

The important qualitative points are these:
\begin{enumerate}[leftmargin=*,itemsep=0.25em]
\item The proposal introduces a genuine order-raising operation; it is not only
      a larger decimal constant.
\item The final height is huge but constructive and independent of the named
      incumbent value.
\item The proof never assumes monotonicity of reflection in its oracle.  That
      monotonicity is false in general because reflected programs may inspect
      oracle values.
\item The huge height remains symbolic in the comparison proof.  Kernel
      checking should not evaluate \(\Pfive\).
\end{enumerate}

%=============================================================================
\section{Repository baseline and constraints}

\subsection{Audited revision}

The development targets \repo{} commit
\code{86b360648163f0249eab79539438b6a83dfe8215}.  The repository pins
\code{leanprover/lean4:v4.34.0}.  At that revision the relevant modules are:
\begin{center}
\begin{tabular}{>{\ttfamily}p{0.32\linewidth}p{0.58\linewidth}}
\toprule
Reflection.lean & Defines oracle evaluation, finite-prefix reflection, and
                  W-tree scheduling.\\
ReflectionFacts.lean & Proves visible oracle values, visible programs, seed
                       values, and child stages are strictly exceeded.\\
Contender6.lean & Defines the universal tree index, the P5 seed family, the
                  final family, and the incumbent.\\
Comparison.lean & Proves the incumbent exceeds P5 and the exact Contender~5
                  port, plus structural domination results.\\
\bottomrule
\end{tabular}
\end{center}

The incumbent definitions reduce to:
\begin{lstlisting}
def p5Family : Family :=
  fun _ n => BigNumber.iterate BigNumber.diagonal n n

def finalFamily (n : Nat) : Nat :=
  stage p5Family index n

def contender_6 : Nat := finalFamily 42
\end{lstlisting}
The universal index is
\begin{lstlisting}
def treesOfDepth (k : Nat) : WTree Nat :=
  .node fun i => defTree (nthD (BigNumber.allTerms k) i .zero)

def index : WTree Nat := .node treesOfDepth
\end{lstlisting}
The first tree under the first depth join is definitionally a leaf:
\begin{lstlisting}
theorem first_tree_leaf :
  defTree (nthD (BigNumber.allTerms 0) 0 .zero) = .leaf := rfl
\end{lstlisting}
This two-edge spine,
\[
  \mathrm{leaf}\longrightarrow \mathrm{treesOfDepth}(0)
  \longrightarrow \Idx,
\]
will be reused at every scheduler order.

\subsection{Contest constraints relevant to the proposal}

The proposal is designed around the following repository requirements.
\begin{enumerate}[leftmargin=*]
\item The number must be a closed, computable \lean{Nat}.
\item No partial, unsafe, external, noncomputable, or oracle stand-in is
      permitted.
\item A theorem must strictly compare the new number with the incumbent.
\item The defining closure must not smuggle in the incumbent value as a huge
      iteration counter.
\item Definition and proof checking must remain symbolic enough to fit the
      contest's build limits.
\item The formal argument must respect the repository's constructive axiom
      policy.
\end{enumerate}

The construction below was selected partly because it addresses all six at the
level of its architecture.  In particular, the comparison theorem unfolds only
a successor constructor around the height; it never asks Lean to compute the
height.

%=============================================================================
\section{The incumbent reflector}

\subsection{Families and visible maxima}

Write
\[
 \Fam=\Nat\to\Nat\to\Nat
\]
for the type called \lean{Family}.  For \(a\in\Fam\), the repository's
\lean{reflect} function at input \(n\) includes the finite prefix
\[
 \max_{i\le n} a(i,n)
\]
and the maximum value of every enumerated object-language term of depth at most
\(n\), evaluated with \(a\) as an oracle family.  It then adds \(n+1\).
Consequently the repository proves
\begin{equation}\label{eq:oracle}
 i\le n\quad\Longrightarrow\quad a(i,n)<\operatorname{reflect}(a,n).
\end{equation}

The crucial point is that Equation~\eqref{eq:oracle} is a direct maximum
inequality.  It does not require \lean{reflect} to be monotone in \(a\).  Such
monotonicity would be unsafe to assume: a program can compare or otherwise
inspect oracle values and follow a different branch when the oracle changes.

\subsection{Forest scheduling}

For a W-tree \(t\in\Tree=\operatorname{WTree}(\Nat)\), the incumbent scheduler
is
\[
 \operatorname{stage}(a,\mathrm{leaf},n)=\operatorname{reflect}(a,n),
\]
\[
 \operatorname{stage}(a,\mathrm{node}(f),n)
 =\operatorname{reflect}\bigl((i,m)\mapsto
   \operatorname{stage}(a,f(i),m),n\bigr).
\]
Equation~\eqref{eq:oracle} immediately gives two facts already proved in
\lean{ReflectionFacts.lean}:
\begin{align}
 a(i,n)&<\operatorname{stage}(a,\mathrm{leaf},n)
 &&(i\le n),\label{eq:seedleaf}\\
 \operatorname{stage}(a,f(i),n)
 &<\operatorname{stage}(a,\mathrm{node}(f),n)
 &&(i\le n).\label{eq:childparent}
\end{align}

Applying these along the two-edge spine inside \(\Idx\) proves that the
incumbent final family exceeds the P5 seed family.  The new proposal repeats the
same local reasoning one type level higher, with schedulers replacing ordinary
stages.

%=============================================================================
\section{Scheduler-level reflection}

\subsection{The scheduler type}

A scheduler has type
\[
 \mathsf{Scheduler}
 =\Fam\to\Tree\to\Nat\to\Nat.
\]
This is exactly the Lean abbreviation:
\begin{lstlisting}
abbrev StageOperator := Family → WTree Nat → Nat → Nat
\end{lstlisting}
The incumbent \lean{stage} is one such scheduler.  The second argument need not
be interpreted only as an object-level ordinal control.  It can also index a
family of recursively promoted schedulers.

\subsection{One order lift}

\begin{definition}[Scheduler lift]\label{def:lift}
For a scheduler \(S\), seed family \(a\), W-tree \(u\), and input \(n\), define
\[
 (\Lift S)(a,\mathrm{leaf},n)=S(a,\Idx,n),
\]
and
\[
 (\Lift S)(a,\mathrm{node}(f),n)
 =S\bigl((i,m)\mapsto(\Lift S)(a,f(i),m),\Idx,n\bigr).
\]
\end{definition}

The Lean definition is almost a transliteration:
\begin{lstlisting}
def liftStage (S : StageOperator) (seed : Family) :
    WTree Nat → Nat → Nat
  | .leaf => S seed index
  | .node f =>
      S (fun i => liftStage S seed (f i)) index
\end{lstlisting}
The recursive call is on the immediate W-subtree \(f(i)\), just as in the
repository's existing \lean{stage}.  There is no general recursion and no
termination proof beyond the W-tree recursor.

\begin{remark}
The second W-tree in \(\Lift S\) controls \emph{operator nesting}.  At a leaf,
we run the entire previous scheduler through the universal object-level index.
At a node, the recursively promoted child operators become the oracle family
seen by the previous scheduler.  Thus one lift is not merely one more object
stage; it reflects the scheduler's complete behavior.
\end{remark}

\subsection{Finite orders}

\begin{definition}[Reflection-power hierarchy]\label{def:orders}
Let
\[
 S_0=\operatorname{stage},\qquad S_{k+1}=\Lift S_k.
\]
\end{definition}
In Lean:
\begin{lstlisting}
def stageOrder : Nat → StageOperator
  | 0 => stage
  | k + 1 => liftStage (stageOrder k)
\end{lstlisting}
This is primitive recursion on \lean{Nat}.  For every concrete natural \(k\),
\(S_k\) is a total scheduler.

%=============================================================================
\section{The proposed number}

\subsection{A seeded height}

Let
\[
 \Pfive=\operatorname{p5Family}(0,42)
 =\operatorname{iterate}(\operatorname{diagonal},42,42).
\]
This is definitionally the numerical expression used by P5's proposal, but the
new definition refers to the seed family rather than to a named contender
constant.  Define
\begin{equation}\label{eq:height}
 \Height=\operatorname{succ}(42+\Pfive).
\end{equation}
The successor wrapper has two purposes.  It gives a large positive height and,
more importantly for the comparison proof, exposes positivity by one
constructor without normalizing \(\Pfive\).

The exact Lean definition is:
\begin{lstlisting}
def metaHeight : Nat := Nat.succ (42 + p5Family 0 42)
\end{lstlisting}

\subsection{Contender 7}

\begin{construction}[Seeded reflection power]\label{con:c7}
Define
\[
 C_7=S_{\Height}(\operatorname{p5Family},\Idx,42).
\]
\end{construction}
Equivalently:
\begin{lstlisting}
def contender_7 : Nat :=
  stageOrder metaHeight p5Family index 42
\end{lstlisting}

This definition is intentionally compact.  The large number comes from the
normalization behavior of a small total program, not from a large source file or
an embedded numeral.

\subsection{Why the number is computable in principle}

\begin{proposition}\label{prop:computable}
\(C_7\) is a well-defined computable natural number.
\end{proposition}
\begin{proof}
The P5 interpreter, term enumeration, finite maximum, diagonal function,
\lean{reflect}, and \lean{stage} are total definitions already present in the
repository.  The new function \lean{liftStage} is structural recursion on an
inductive W-tree.  The function \lean{stageOrder} is primitive recursion on a
natural number.  The height \(\Height\) is a closed application of total
functions, and every remaining argument in Construction~\ref{con:c7} is closed.
Therefore normalization terminates in the metatheory, although the required
resources are incomprehensibly large.
\end{proof}

\begin{caution}
``Computable in principle'' does not mean feasible.  Evaluating \(C_7\), or even
materializing more than tiny toy instances of the hierarchy, is not a practical
objective.  The contest compares closed normal forms by proof, not by execution.
\end{caution}

%=============================================================================
\section{Formal comparison architecture}

This section gives the complete mathematical proof mirrored by
\lean{Comparison7.lean}.

\subsection{The domination invariant}

\begin{definition}\label{def:dom}
A scheduler \(S\) is \emph{index-dominating} if
\[
 \operatorname{Dom}(S):\iff
 \forall a\in\Fam\ \forall i,n\in\Nat,
 \quad i\le n\Longrightarrow a(i,n)<S(a,\Idx,n).
\]
\end{definition}
The Lean predicate is:
\begin{lstlisting}
def IndexDominates (S : StageOperator) : Prop :=
  ∀ (seed : Family) {i n : Nat},
    i ≤ n → seed i n < S seed index n
\end{lstlisting}

\subsection{Base scheduler}

\begin{lemma}\label{lem:base}
The incumbent scheduler \(S_0=\operatorname{stage}\) is index-dominating.
\end{lemma}
\begin{proof}
Fix \(a,i,n\) with \(i\le n\).  By Equation~\eqref{eq:seedleaf},
\[
 a(i,n)<\operatorname{stage}(a,\mathrm{leaf},n).
\]
The leaf is the zeroth child of \(\operatorname{treesOfDepth}(0)\), and the
latter is the zeroth child of \(\Idx\).  Applying
Equation~\eqref{eq:childparent} twice gives
\[
 \operatorname{stage}(a,\mathrm{leaf},n)
 <\operatorname{stage}(a,\operatorname{treesOfDepth}(0),n)
 <\operatorname{stage}(a,\Idx,n).
\]
Transitivity proves the claim.
\end{proof}

The Lean proof reuses \lean{first_tree_leaf} exactly as the incumbent comparison
proof does.  It is generic in the seed family, which is what makes the
higher-order induction possible.

\subsection{A child is below its promoted parent}

\begin{lemma}\label{lem:child}
Assume \(\operatorname{Dom}(S)\).  If \(i\le n\), then
\[
 (\Lift S)(a,f(i),n)<(\Lift S)(a,\mathrm{node}(f),n).
\]
\end{lemma}
\begin{proof}
Unfold the node clause of Definition~\ref{def:lift}.  Put
\[
 b(j,m)=(\Lift S)(a,f(j),m).
\]
The right side becomes \(S(b,\Idx,n)\), while the left side is \(b(i,n)\).
The claim is precisely \(\operatorname{Dom}(S)\) applied to \(b\) and
\(i\le n\).
\end{proof}

This is the key higher-order step.  It uses only the domination invariant; it
does not compare two unrelated oracle families.

\subsection{Leaf lift below universal-index lift}

\begin{lemma}\label{lem:leafindex}
If \(\operatorname{Dom}(S)\), then for every \(a,n\),
\[
 (\Lift S)(a,\mathrm{leaf},n)<(\Lift S)(a,\Idx,n).
\]
\end{lemma}
\begin{proof}
Apply Lemma~\ref{lem:child} along the same two-edge spine used in
Lemma~\ref{lem:base}:
\[
 (\Lift S)(a,\mathrm{leaf},n)
 <(\Lift S)(a,\operatorname{treesOfDepth}(0),n)
 <(\Lift S)(a,\Idx,n).
\]
\end{proof}

\subsection{The invariant is preserved}

\begin{lemma}\label{lem:preserve}
If \(\operatorname{Dom}(S)\), then \(\operatorname{Dom}(\Lift S)\).
\end{lemma}
\begin{proof}
For \(i\le n\), first use \(\operatorname{Dom}(S)\):
\[
 a(i,n)<S(a,\Idx,n).
\]
By the leaf equation in Definition~\ref{def:lift},
\[
 S(a,\Idx,n)=(\Lift S)(a,\mathrm{leaf},n).
\]
Lemma~\ref{lem:leafindex} then gives
\[
 a(i,n)<(\Lift S)(a,\Idx,n).
\]
\end{proof}

\begin{corollary}\label{cor:allorders}
Every scheduler order \(S_k\) is index-dominating.
\end{corollary}
\begin{proof}
Induction on \(k\), using Lemma~\ref{lem:base} at zero and
Lemma~\ref{lem:preserve} at successors.
\end{proof}

\subsection{Strict growth at every successor order}

\begin{theorem}\label{thm:step}
For every seed family \(a\), input \(n\), and order \(k\),
\[
 S_k(a,\Idx,n)<S_{k+1}(a,\Idx,n).
\]
\end{theorem}
\begin{proof}
By the leaf equation,
\[
 S_k(a,\Idx,n)=(\Lift S_k)(a,\mathrm{leaf},n).
\]
Corollary~\ref{cor:allorders} gives \(\operatorname{Dom}(S_k)\), so
Lemma~\ref{lem:leafindex} yields
\[
 (\Lift S_k)(a,\mathrm{leaf},n)
 <(\Lift S_k)(a,\Idx,n)=S_{k+1}(a,\Idx,n).
\]
\end{proof}

\begin{corollary}\label{cor:positive}
For every \(h>0\),
\[
 S_0(a,\Idx,n)<S_h(a,\Idx,n).
\]
\end{corollary}
The Lean development packages this as a recursion over a predecessor:
\begin{lstlisting}
theorem stageOrder_zero_lt_succ (seed : Family) (n : Nat) :
    ∀ k,
      stageOrder 0 seed index n <
      stageOrder (Nat.succ k) seed index n
\end{lstlisting}

\subsection{The required comparison}

\begin{theorem}[Proposed kernel theorem]\label{thm:main}
\[
 \texttt{BusyLean.contender\_6}
 <\texttt{BusyLean.contender\_7}.
\]
\end{theorem}
\begin{proof}
The incumbent is definitionally
\[
 C_6=S_0(\operatorname{p5Family},\Idx,42).
\]
The proposed number is
\[
 C_7=S_{\operatorname{succ}(42+\Pfive)}
       (\operatorname{p5Family},\Idx,42).
\]
Apply Corollary~\ref{cor:positive} with predecessor \(42+\Pfive\).
\end{proof}

The Lean proof is correspondingly short:
\begin{lstlisting}
theorem contender_6_lt_contender_7 :
    contender_6 < contender_7 := by
  change stageOrder 0 p5Family index 42 <
    stageOrder (Nat.succ (42 + p5Family 0 42))
      p5Family index 42
  exact stageOrder_zero_lt_succ
    p5Family 42 (42 + p5Family 0 42)
\end{lstlisting}
No reduction of \lean{p5Family 0 42} is required.  The term appears only as a
symbolic argument to a theorem quantified over every natural predecessor.

\subsection{Inherited comparisons}

Once Theorem~\ref{thm:main} is accepted, the existing repository theorems give
immediately:
\[
 \texttt{contender\_5}<\texttt{contender\_7}
\]
and
\[
 \texttt{BigNumber.contender\_6}<\texttt{contender\_7}.
\]
The latter name is P5's proposal-local constant, not the BusyLean incumbent.
The archive includes both transitivity proofs.

%=============================================================================
\section{Magnitude analysis}

\subsection{What is formal}

Only the following magnitude claims are intended as formal Lean theorems:
\begin{enumerate}[leftmargin=*]
\item every scheduler order is index-dominating;
\item each successor scheduler order is strictly larger at the universal index;
\item the incumbent is strictly below the proposed number;
\item the exact Contender~5 port and P5's proposal-local number are also below
      the proposed number by transitivity.
\end{enumerate}
These are ordinary inequalities in \(\Nat\).  They do not require ordinal
notations, well-ordering proofs, or a formal fast-growing hierarchy.

\subsection{Repository calibration}

For scale only, adopt the current repository article's notation
\[
 \Lambda\approx\psi(\Omega_\omega)
\]
for the strength of P5's W-language.  The repository estimates
\[
 \Pfive\quad\hbox{near}\quad f_{\Lambda+\omega^2}(42)
\]
and the incumbent near
\[
 C_6\quad\hbox{near}\quad f_{\Lambda\cdot\Lambda}(42)
 =f_{\Lambda^2}(42).
\]
These are calibrations, not proved equalities.

\subsection{Effect of one scheduler lift}

A stage along \(\Idx\) traverses an index whose available W-ranks are calibrated
by \(\Lambda\).  A scheduler lift traverses a fresh copy of that index while
each node invokes the complete preceding scheduler.  The natural ordinal
bookkeeping model is therefore right multiplication by \(\Lambda\):
\[
 \alpha\longmapsto\alpha\cdot\Lambda.
\]
Starting from \(\Lambda^2\), this suggests
\[
 S_h(\operatorname{p5Family},\Idx,42)
 \quad\hbox{near}\quad f_{\Lambda^{h+2}}(42)
\]
for finite \(h\).  At the seeded height \(\Height\), the proposal is therefore
heuristically
\begin{equation}\label{eq:heuristic}
 C_7\quad\hbox{near}\quad f_{\Lambda^{\Height+2}}(42),
 \qquad
 \Height=\operatorname{succ}(42+\Pfive).
\end{equation}

Even a fixed height of~42 would move the calibration from \(\Lambda^2\) to
roughly \(\Lambda^{44}\).  The chosen height is itself the P5 seed value plus a
guard, producing an incomparably larger finite exponent.

\begin{caution}
Equation~\eqref{eq:heuristic} is intentionally not advertised as a theorem.  A
rigorous calibration would require a formal ordinal semantics for the object
language, the universal index, oracle interpretation, and scheduler lift.  The
finite exponent \(\Height\) is an actual natural number, however enormous; the
claim is not that the construction reaches a literal transfinite iteration
count.
\end{caution}

\subsection{Why this is more than ``use the old number as a counter''}

The construction has two independent sources of strength.
\begin{enumerate}[leftmargin=*]
\item \textbf{New operation.} \(S\mapsto\Lift S\) raises the order of reflection
      by recursively turning promoted child schedulers into the previous
      scheduler's oracle family.  This operation already yields an infinite
      strict hierarchy when iterated through ordinary small naturals.
\item \textbf{Seeded scale.} The final order count is synchronized with the
      inherited P5 diagonal value.  It is not the BusyLean incumbent and is not
      needed to establish the hierarchy's strictness; it only selects a very
      high member of the new hierarchy.
\end{enumerate}
A conservative variant \(S_{42}(\operatorname{p5Family},\Idx,42)\) is therefore
available if a contest judge prefers a fixed promotion count.  The main package
uses the seeded height because the repository rules prohibit the incumbent
value specifically, not all previously defined constructive growth functions,
and because the strictness proof remains entirely symbolic.

%=============================================================================
\section{Engineering analysis}

\subsection{Expected definition-check behavior}

The defining module contains only an abbreviation and three short recursive
definitions.  Lean checks the structural equations for \lean{liftStage} and the
primitive equations for \lean{stageOrder}; it need not normalize
\lean{metaHeight} or \lean{contender_7}.  The source therefore has a plausible
path to the repository's definition-time limit.

\subsection{Expected proof-check behavior}

The proof introduces a generic invariant and proves it once.  The final theorem
applies a universally quantified lemma at the symbolic predecessor
\lean{42 + p5Family 0 42}.  It does not construct a transitivity chain of that
length and does not reduce the P5 term.  Thus proof size is constant in the
numerical value of the height.

\subsection{Axiom and dependency audit}

The supplied audit computes the constant closure of
\lean{BusyLean.contender_7}.  It is intended to fail if that closure includes
any of:
\begin{lstlisting}
BusyLean.contender_6
BusyLean.contender_5
BigNumber.contender_6
BigNumber.contender_5
\end{lstlisting}
Reusing \lean{p5Family}, \lean{index}, and the total language implementation is
allowed and visible.  Reusing a named predecessor value as an opaque subterm is
not.

The audit also runs \texttt{\#print axioms} on the new number and comparison
theorems.  The expected axiom profile is inherited from the accepted upstream
modules; the proposal adds no axiom declaration, choice operator, or external
implementation hook.

\subsection{Why no monotonicity lemma appears}

It would be tempting to prove that a larger oracle produces a larger reflected
value and then lift inequalities pointwise.  The upstream source explicitly
warns that this principle fails: interpreted programs may test oracle values.
The present proof instead arranges for each desired child value to be a
\emph{visible component of the exact oracle family} passed to the parent.
Equation~\eqref{eq:oracle} then gives a strict inequality directly.  This is
both more constructive and more robust.

\subsection{Integration footprint}

The overlay adds:
\begin{itemize}[leftmargin=*]
\item \code{BusyLean/Contender7.lean};
\item \code{BusyLean/Comparison7.lean};
\item two imports in \code{BusyLean.lean};
\item \code{test/Contender7Audit.lean}.
\end{itemize}
No upstream file needs semantic modification.  The archive contains a complete
Git patch against the audited revision.

%=============================================================================
\section{Independent executable checks}

Because a Lean executable was unavailable, the package includes checks that are
useful but explicitly weaker than kernel verification.

\subsection{Static source audit}

The static audit verifies:
\begin{enumerate}[leftmargin=*]
\item absence of \code{noncomputable}, \code{partial}, \code{unsafe},
      \code{extern}, \code{implemented_by}, \code{native_decide},
      \code{sorry}, \code{admit}, and custom \code{axiom} declarations from the
      production overlay;
\item absence of the token \code{contender_6} from the defining module;
\item presence of the intended defining equations and all comparison lemmas;
\item balanced delimiters and the kernel closure test.
\end{enumerate}
The check passes for the files in this archive.

\subsection{Finite reflection-power model}

A memoized Python model replaces infinitely branching W-trees by finite trees
with a leaf default and retains the exact recursive architecture of
\lean{stageOrder}.  Its reflector exceeds every visible oracle value, matching
the only property used by the formal proof.  It tests scheduler orders
\(0,1,2,3\) and inputs \(0,1,2\).

\begin{center}
\begin{tabular}{crrrr}
\toprule
Input & Order 0 & Order 1 & Order 2 & Order 3\\
\midrule
0 & 6 & 12 & 30 & 84\\
1 & 14 & 38 & 134 & 518\\
2 & 24 & 60 & 204 & 780\\
\bottomrule
\end{tabular}
\end{center}
Every row is strictly increasing.  The model also checks, at every tested input
and order, the exact proof spine
\[
 (\Lift S)(a,\mathrm{leaf},n)
 <(\Lift S)(a,\operatorname{treesOfDepth}(0),n)
 <(\Lift S)(a,\Idx,n).
\]
This catches errors in recursion direction and oracle wiring, but it does not
parse Lean and is not evidence about kernel acceptance.

\subsection{Pinned CI workflow}

The archive provides a GitHub Actions workflow that uses the repository's
\code{lean-toolchain}, deletes \code{.lake}, runs \code{lake build}, and then
runs the closure/axiom audit.  This is the intended final acceptance test.

%=============================================================================
\section{Risks and unresolved verification points}

A serious proposal should distinguish mathematical risk from environment risk.

\subsection{Lean elaboration risk}

The source was written by matching the exact signatures of
\lean{stage_child_lt}, \lean{seed_lt_stage_leaf},
\lean{first_tree_leaf}, \lean{treesOfDepth}, and \lean{index} at the audited
commit.  Nevertheless, only the Lean elaborator and kernel can confirm all
implicit arguments, definitional reductions, and equation-compiler behavior.
The supplied code is compact enough that any elaboration repair should be local,
but no repair is silently assumed here.

\subsection{Timing risk}

The design is intended to avoid normalization of huge terms, but the contest
limits must be measured on the pinned toolchain from a clean build.  The CI
workflow should record wall-clock times for the defining module, comparison
module, and audit separately before repository acceptance.

\subsection{Calibration risk}

The ordinal estimate depends on the repository's informal calibration and on
the modeling choice \(\alpha\mapsto\alpha\Lambda\) for one scheduler lift.  A
formal semantics might produce a different expression.  This does not affect
the strict natural-number comparison.

\subsection{Judging-policy risk}

A judge might prefer a fixed order such as~42 rather than a P5-scaled order,
even though the stated anti-laziness rule bars the incumbent value rather than
all large constructive counters.  The new operator and all proofs remain valid
under the fixed-height variant.  The archive documents this fallback as a
policy variant, not as the primary proposal.

%=============================================================================
\section{Further research questions}

\begin{enumerate}[leftmargin=*,label=\textbf{R\arabic*.}]
\item \textbf{Kernel closure and timing.} Does the overlay build unchanged on
      Lean~4.34.0, and what are the clean-build times of the definition,
      comparison, and audit modules?

\item \textbf{Formal ordinal semantics.} Construct an ordinal notation system
      and interpretation for the P5 term language, \lean{stage}, the universal
      index, and \lean{liftStage}.  Prove sound upper and lower bounds for each
      scheduler order rather than relying on the heuristic
      \(\Lambda^{h+2}\).

\item \textbf{Fixed-height versus seeded-height policy.} Determine a contest
      standard for large iteration counts inherited from earlier constructive
      languages.  A formal ``non-laziness margin'' predicate could make this
      judgment objective.

\item \textbf{Operator-level universal trees.} The present hierarchy iterates
      the same lift by a natural count.  Can one define a conservative language
      of scheduler expressions whose own W-trees index heterogeneous sequences
      of order-raising operators, while retaining an internal normalization
      proof?

\item \textbf{Transfinite promotion without nontermination.} Develop a notation
      system for well-founded trees of scheduler lifts and an evaluator by
      structural recursion.  This could replace the finite count \(\Height\)
      with a directly represented well-founded control tree.

\item \textbf{System F and higher universes.} The upstream roadmap proposes
      System~F, System~F$_\omega$, and the Calculus of Constructions as
      qualitatively stronger base languages.  How does scheduler reflection
      interact with each normalization proof?  Does ``stronger language first,
      then reflection power'' dominate pure reflection iteration decisively?

\item \textbf{Parametric comparison theorem.} Abstract the two-edge spine of
      \lean{index} into a structure containing a distinguished leaf descendant.
      Prove the entire scheduler-power theorem for any such index and any base
      scheduler satisfying visible-oracle domination.

\item \textbf{Proof-term minimization.} Can the comparison be expressed using a
      reusable algebra of progressive schedulers, reducing both source size and
      kernel-checking time while preserving a transparent closure audit?

\item \textbf{Machine-checked resource bounds.} Formalize a syntactic cost model
      for elaboration and kernel checking, then prove that the comparison proof
      avoids evaluating \lean{p5Family 0 42}.

\item \textbf{Independent implementation.} Reimplement the scheduler-power
      construction in a minimal standalone Lean file and prove agreement with
      the repository overlay, analogous to the existing port-agreement audit.
\end{enumerate}

%=============================================================================
\section{Conclusion}

The universal W-tree index in Contender~6 does more than support one forest
reflection.  Because the scheduler itself has the type of a family-to-tree
transformer, the same index can promote the scheduler as an object.  The
resulting lift preserves a simple domination invariant and is strictly
progressive at the universal index.  Primitive recursion therefore gives a
strict hierarchy whose zero level is exactly the incumbent scheduler.

Selecting the P5 seed value, plus a successor guard, as the hierarchy height
produces a closed constructive number with a very large formal margin:
\[
 C_6<C_7.
\]
The proof is structurally small and does not normalize the huge height.  The
proposal is consequently a plausible next BusyLean contender even before the
more ambitious language-strength jumps in the repository roadmap.

The remaining decisive step is mechanical rather than conceptual: run the
provided overlay and audit under the exact Lean~4.34.0 toolchain, record timing,
and accept or locally repair any elaboration issue.  Until that is done, the
archive should be treated as a rigorous research proposal with executable
supporting evidence, not as a completed kernel-certified submission.

%=============================================================================
\appendix
\section{Complete defining module}
\lstinputlisting{../lean/BusyLean/Contender7.lean}

\section{Complete comparison module}
\lstinputlisting{../lean/BusyLean/Comparison7.lean}

\section{Kernel audit module}
\lstinputlisting{../lean/test/Contender7Audit.lean}

\section{Reproduction commands}

From a checkout of the audited BusyLean commit:
\begin{lstlisting}[language={}]
cp lean/BusyLean/Contender7.lean BusyLean/
cp lean/BusyLean/Comparison7.lean BusyLean/
cat lean/BusyLean.lean.append >> BusyLean.lean
cp lean/test/Contender7Audit.lean test/
rm -rf .lake
lake build
lake env lean test/Contender7Audit.lean
\end{lstlisting}
The package-side checks are:
\begin{lstlisting}[language={}]
./scripts/run_checks.sh
\end{lstlisting}

\begin{thebibliography}{9}
\bibitem{BusyLean}
V. Reshetnikov et al.,
\emph{BusyLean: Naming Ever Larger Numbers in Lean}, repository and article,
commit \code{86b360648163f0249eab79539438b6a83dfe8215}, 2026.
\url{https://github.com/VladimirReshetnikov/BusyLean}

\bibitem{Roux}
C. Roux,
\emph{name-the-biggest-number}, Coq large-number contest repository.
\url{https://github.com/codyroux/name-the-biggest-number}

\bibitem{Lean}
The Lean developers,
\emph{Lean 4 theorem prover}, version 4.34.0, 2026.
\url{https://github.com/leanprover/lean4}

\bibitem{SchwichtenbergWainer}
H. Schwichtenberg and S. S. Wainer,
\emph{Proofs and Computations}, Cambridge University Press, 2012.

\bibitem{Girard}
J.-Y. Girard, Y. Lafont, and P. Taylor,
\emph{Proofs and Types}, Cambridge University Press, 1989.

\bibitem{Buchholz}
W. Buchholz,
A new system of proof-theoretic ordinal functions,
\emph{Annals of Pure and Applied Logic} 32 (1986), 195--207.
\end{thebibliography}

\end{document}
